\chapter{$\lambda$-矩阵}

\section{习题}

\begin{enumerate}
  \item 化下列 $\lambda$-矩阵成标准形：
  \begin{enumerate}
    \item \[
    \begin{pmatrix}
      \lambda^3-\lambda & 2\lambda^2 \\
      \lambda^2+5\lambda & 3\lambda
    \end{pmatrix}\text{；}
    \]
  \anssub
    \item \[
    \begin{pmatrix}
      1-\lambda & \lambda^2 & \lambda \\
      \lambda & \lambda & -\lambda \\
      1+\lambda^2 & \lambda^2 & -\lambda^2
    \end{pmatrix}\text{；}
    \]
  \anssub
    \item \[
    \begin{pmatrix}
      \lambda^2+\lambda & 0 & 0 \\
      0 & \lambda & 0 \\
      0 & 0 & (\lambda+1)^2
    \end{pmatrix}\text{；}
    \]
  \anssub
    \item \[
    \begin{pmatrix}
      0 & 0 & 0 & \lambda^2 \\
      0 & 0 & \lambda^2-\lambda & 0 \\
      0 & (\lambda-1)^2 & 0 & 0 \\
      \lambda^2-\lambda & 0 & 0 & 0
    \end{pmatrix}\text{；}
    \]
  \anssub
    \item \[
    \begin{pmatrix}
      3\lambda^2+2\lambda-3 & 2\lambda-1 & \lambda^2+2\lambda-3 \\
      4\lambda^2+3\lambda-5 & 3\lambda-2 & \lambda^2+3\lambda-4 \\
      \lambda^2+\lambda-4 & \lambda-2 & \lambda-1
    \end{pmatrix}\text{；}
    \]
  \anssub
    \item \[
    \begin{pmatrix}
      2\lambda & 3 & 0 & 1 & \lambda \\
      4\lambda & 3\lambda+6 & 0 & \lambda+2 & 2\lambda \\
      0 & 6\lambda & \lambda & 2\lambda & 0 \\
      \lambda-1 & 0 & \lambda-1 & 0 & 0 \\
      3\lambda-3 & 1-\lambda & 2\lambda-2 & 0 & 0
    \end{pmatrix}\text{．}
    \]
  \anssub
  \end{enumerate}

  \item 求下列 $\lambda$-矩阵的不变因子：
  \begin{enumerate}
    \item \[
    \begin{pmatrix}
      \lambda-2 & -1 & 0 \\
      0 & \lambda-2 & -1 \\
      0 & 0 & \lambda-2
    \end{pmatrix}\text{；}
    \]
  \anssub
    \item \[
    \begin{pmatrix}
      \lambda & -1 & 0 & 0 \\
      0 & \lambda & -1 & 0 \\
      0 & 0 & \lambda & -1 \\
      5 & 4 & 3 & \lambda+2
    \end{pmatrix}\text{；}
    \]
  \anssub
    \item \[
    \begin{pmatrix}
      \lambda+\alpha & \beta & 1 & 0 \\
      -\beta & \lambda+\alpha & 0 & 1 \\
      0 & 0 & \lambda+\alpha & \beta \\
      0 & 0 & -\beta & \lambda+\alpha
    \end{pmatrix}\text{；}
    \]
  \anssub
    \item \[
    \begin{pmatrix}
      0 & 0 & 1 & \lambda+2 \\
      0 & 1 & \lambda+2 & 0 \\
      1 & \lambda+2 & 0 & 0 \\
      \lambda+2 & 0 & 0 & 0
    \end{pmatrix}\text{；}
    \]
  \anssub
    \item \[
    \begin{pmatrix}
      \lambda+1 & 0 & 0 & 0 \\
      0 & \lambda+2 & 0 & 0 \\
      0 & 0 & \lambda-1 & 0 \\
      0 & 0 & 0 & \lambda-2
    \end{pmatrix}\text{．}
    \]
  \anssub
  \end{enumerate}

  \item 证明：
  \[
  \begin{pmatrix}
    \lambda & 0 & 0 & \cdots & 0 & a_n \\
    -1 & \lambda & 0 & \cdots & 0 & a_{n-1} \\
    0 & -1 & \lambda & \cdots & 0 & a_{n-2} \\
    \vdots & \vdots & \vdots & & \vdots & \vdots \\
    0 & 0 & 0 & \cdots & \lambda & a_2 \\
    0 & 0 & 0 & \cdots & -1 & \lambda+a_1
  \end{pmatrix}
  \]
  的不变因子是 $\underbrace{1,1,\cdots,1}_{n-1\text{个}},f(\lambda)$，其中 $f(\lambda)=\lambda^n+a_1\lambda^{n-1}+\cdots+a_{n-1}\lambda+a_n$．

  \ansproof
  \item 设 $A$ 是数域 $P$ 上一个 $n\times n$ 矩阵．证明：$A$ 与 $A^{\mathrm{T}}$ 相似．

  \ansproof
  \item 设
  \[
  A=\begin{pmatrix}
    \lambda & 0 & 0 \\
    1 & \lambda & 0 \\
    0 & 1 & \lambda
  \end{pmatrix},
  \]
  求 $A^k$．

  \anscompute
  \item 求下列复矩阵的若尔当标准形：
  \begin{enumerate}
    \item \[
    \begin{pmatrix}
      1 & 2 & 0 \\
      0 & 2 & 0 \\
      -2 & -2 & -1
    \end{pmatrix}\text{；}
    \]
  \anssub
    \item \[
    \begin{pmatrix}
      13 & 16 & 16 \\
      -5 & -7 & -6 \\
      -6 & -8 & -7
    \end{pmatrix}\text{；}
    \]
  \anssub
    \item \[
    \begin{pmatrix}
      3 & 0 & 8 \\
      3 & -1 & 6 \\
      -2 & 0 & -5
    \end{pmatrix}\text{；}
    \]
  \anssub
    \item \[
    \begin{pmatrix}
      4 & 5 & -2 \\
      -2 & -2 & 1 \\
      -1 & -1 & 1
    \end{pmatrix}\text{；}
    \]
  \anssub
    \item \[
    \begin{pmatrix}
      3 & 7 & -3 \\
      -2 & -5 & 2 \\
      -4 & -10 & 3
    \end{pmatrix}\text{；}
    \]
  \anssub
    \item \[
    \begin{pmatrix}
      1 & -1 & 2 \\
      3 & -3 & 6 \\
      2 & -2 & 4
    \end{pmatrix}\text{；}
    \]
  \anssub
    \item \[
    \begin{pmatrix}
      1 & 1 & -1 \\
      -3 & -3 & 3 \\
      -2 & -2 & 2
    \end{pmatrix}\text{；}
    \]
  \anssub
    \item \[
    \begin{pmatrix}
      -4 & 2 & 10 \\
      -4 & 3 & 7 \\
      -3 & 1 & 7
    \end{pmatrix}\text{；}
    \]
  \anssub
    \item \[
    \begin{pmatrix}
      0 & 3 & 3 \\
      -1 & 8 & 6 \\
      2 & -14 & -10
    \end{pmatrix}\text{；}
    \]
  \anssub
    \item \[
    \begin{pmatrix}
      8 & 30 & -14 \\
      -6 & -19 & 9 \\
      -6 & -23 & 11
    \end{pmatrix}\text{；}
    \]
  \anssub
    \item \[
    \begin{pmatrix}
      3 & 1 & 0 & 0 \\
      -4 & -1 & 0 & 0 \\
      7 & 1 & 2 & 1 \\
      -7 & -6 & -1 & 0
    \end{pmatrix}\text{；}
    \]
  \anssub
    \item \[
    \begin{pmatrix}
      1 & 2 & 3 & 4 \\
      0 & 1 & 2 & 3 \\
      0 & 0 & 1 & 2 \\
      0 & 0 & 0 & 1
    \end{pmatrix}\text{；}
    \]
  \anssub
    \item \[
    \begin{pmatrix}
      1 & -3 & 0 & 3 \\
      -2 & 6 & 0 & 13 \\
      0 & -3 & 1 & 3 \\
      -1 & 2 & 0 & 8
    \end{pmatrix}\text{；}
    \]
  \anssub
    \item \[
    \begin{pmatrix}
      0 & 1 & 0 & \cdots & 0 & 0 \\
      0 & 0 & 1 & \cdots & 0 & 0 \\
      \vdots & \vdots & \vdots & & \vdots & \vdots \\
      0 & 0 & 0 & \cdots & 0 & 1 \\
      1 & 0 & 0 & \cdots & 0 & 0
    \end{pmatrix}\text{．}
    \]
  \anssub
  \end{enumerate}

  \item 把习题 6 中各矩阵看成有理数域上矩阵，试写出它们的有理标准形．
  \anscompute
\end{enumerate}

\section{补充题}

\begin{enumerate}
  \item 设 $\mathcal{A}$ 是 $n$ 维线性空间 $V$ 上的线性变换．证明：
  \begin{enumerate}
    \item 若 $\mathcal{A}$ 在 $V$ 的一组基下的矩阵 $A$ 是多项式 $d(\lambda)$ 的友矩阵，则 $\mathcal{A}$ 的最小多项式是 $d(\lambda)$；
  \anssub
    \item 设 $\mathcal{A}$ 的最高次的不变因子是 $d(\lambda)$，则 $\mathcal{A}$ 的最小多项式是 $d(\lambda)$．
  \anssub
  \end{enumerate}
\end{enumerate}
